% Options for packages loaded elsewhere
\PassOptionsToPackage{unicode}{hyperref}
\PassOptionsToPackage{hyphens}{url}
\PassOptionsToPackage{dvipsnames,svgnames,x11names}{xcolor}
\documentclass[
  spanish,
  10pt,
]{article}
\usepackage{xcolor}
\usepackage[margin=2.2cm]{geometry}
\usepackage{amsmath,amssymb}
\setcounter{secnumdepth}{-\maxdimen} % remove section numbering
\usepackage{iftex}
\ifPDFTeX
  \usepackage[T1]{fontenc}
  \usepackage[utf8]{inputenc}
  \usepackage{textcomp} % provide euro and other symbols
\else % if luatex or xetex
  \usepackage{unicode-math} % this also loads fontspec
  \defaultfontfeatures{Scale=MatchLowercase}
  \defaultfontfeatures[\rmfamily]{Ligatures=TeX,Scale=1}
\fi
\usepackage{lmodern}
\ifPDFTeX\else
  % xetex/luatex font selection
\fi
% Use upquote if available, for straight quotes in verbatim environments
\IfFileExists{upquote.sty}{\usepackage{upquote}}{}
\IfFileExists{microtype.sty}{% use microtype if available
  \usepackage[]{microtype}
  \UseMicrotypeSet[protrusion]{basicmath} % disable protrusion for tt fonts
}{}
\makeatletter
\@ifundefined{KOMAClassName}{% if non-KOMA class
  \IfFileExists{parskip.sty}{%
    \usepackage{parskip}
  }{% else
    \setlength{\parindent}{0pt}
    \setlength{\parskip}{6pt plus 2pt minus 1pt}}
}{% if KOMA class
  \KOMAoptions{parskip=half}}
\makeatother
\usepackage{longtable,booktabs,array}
\usepackage{calc} % for calculating minipage widths
% Correct order of tables after \paragraph or \subparagraph
\usepackage{etoolbox}
\makeatletter
\patchcmd\longtable{\par}{\if@noskipsec\mbox{}\fi\par}{}{}
\makeatother
% Allow footnotes in longtable head/foot
\IfFileExists{footnotehyper.sty}{\usepackage{footnotehyper}}{\usepackage{footnote}}
\makesavenoteenv{longtable}
\ifLuaTeX
\usepackage[bidi=basic,shorthands=off,]{babel}
\else
\usepackage[bidi=default,shorthands=off,]{babel}
\fi
\ifLuaTeX
  \usepackage{selnolig} % disable illegal ligatures
\fi
\setlength{\emergencystretch}{3em} % prevent overfull lines
\providecommand{\tightlist}{%
  \setlength{\itemsep}{0pt}\setlength{\parskip}{0pt}}
% Preámbulo para compilar doble_ponderacion.md a PDF con XeTeX (Tectonic).
%
% Resuelve los 312 caracteres que la fuente por defecto (Latin Modern) no tiene
% y que salían como hueco en blanco:
%
%   ─ │ ├ └ ┐ ┘   (144+60+38+18+2+2) el diagrama de la sección 2
%   ✓             (34) la tabla de recuento de la sección 2.1
%   α ρ           (10+4) el parámetro de la sección 5.1 y la rho de Spearman
%
% Dos estrategias distintas a propósito:
%
%  - Los símbolos matemáticos (α ρ ✓) se mapean a comandos de LaTeX, que
%    componen igual en cualquier equipo y no dependen de ninguna fuente
%    instalada.
%  - Los de dibujo de caja NO tienen equivalente en LaTeX, así que hacen falta
%    de una fuente real. Consolas viene con Windows desde Vista y los trae
%    todos; se verificó que existe antes de fijarla aquí.

\usepackage{newunicodechar}
\newunicodechar{α}{\ensuremath{\alpha}}
\newunicodechar{ρ}{\ensuremath{\rho}}
\newunicodechar{✓}{\ensuremath{\checkmark}}

% Consolas sólo para el texto monoespaciado: es donde vive el diagrama. El
% cuerpo del documento sigue con la serif de LaTeX, que compone mucho mejor.
\setmonofont{Consolas}[Scale=0.82]

% El diagrama de la sección 2 mide 74 columnas y se sale de la caja con el
% cuerpo por defecto. Reducirlo aquí evita tener que estrechar el margen de
% todo el documento por culpa de un solo bloque.
\usepackage{fancyvrb}
\fvset{fontsize=\footnotesize}

% Las 18 tablas son anchas; sin esto el texto de las celdas queda con espacios
% enormes entre palabras.
\usepackage{ragged2e}
\usepackage{array}

% Deja partir las palabras en monoespaciada. Sin esto, el nombre del PDF de
% origen --63 caracteres sin espacios-- no cabe en la caja y se sale 18,45 pt
% por el margen derecho, que es visible en la página 1.
\usepackage[htt]{hyphenat}

% Filetes de tabla con aire, que es lo que separa una tabla legible de una
% rejilla apretada.
\usepackage{booktabs}
\setlength{\aboverulesep}{2pt}
\setlength{\belowrulesep}{2pt}
\renewcommand{\arraystretch}{1.15}
\usepackage{bookmark}
\IfFileExists{xurl.sty}{\usepackage{xurl}}{} % add URL line breaks if available
\urlstyle{same}
\hypersetup{
  pdflang={es},
  colorlinks=true,
  linkcolor={RoyalBlue},
  filecolor={Maroon},
  citecolor={Blue},
  urlcolor={Blue},
  pdfcreator={LaTeX via pandoc}}

\author{}
\date{}

\begin{document}

{
\setcounter{tocdepth}{2}
\tableofcontents
}
\section{¿Se están ponderando dos veces los mismos
elementos?}\label{se-estuxe1n-ponderando-dos-veces-los-mismos-elementos}

Revisión de las ponderaciones nuevas del programa de priorización contra
la metodología oficial del índice de Riesgo de Incendios Forestales de
CONAF.

\textbf{Fecha:} 2026-09-14 \textbf{Documento revisado:}
\texttt{Riesgo-Incendios-Forestales-Resumen-Ejecutivo-Final-12022021.pdf}
(Departamento de Desarrollo e Investigación, GEPRIF, CONAF, enero 2021;
37 páginas). Copia en \texttt{INSUMO\_PRIORIZACION/}. \textbf{Código
revisado:} \texttt{lab/priorizacion/notebook/priorizacion.py}, capas de
\texttt{lab/priorizacion/data/raw/} y \texttt{raw2/}, y la capa
publicada \texttt{priorizacion.geojson}.

\begin{center}\rule{0.5\linewidth}{0.5pt}\end{center}

\subsection{1. Respuesta corta}\label{respuesta-corta}

\textbf{Sí, hay doble ponderación, y afecta a dos de los cinco
componentes.} No a los cinco, y no en el mismo grado. El detalle importa
más que el titular:

\begin{longtable}[]{@{}
  >{\raggedright\arraybackslash}p{(\linewidth - 6\tabcolsep) * \real{0.2308}}
  >{\raggedleft\arraybackslash}p{(\linewidth - 6\tabcolsep) * \real{0.3077}}
  >{\raggedright\arraybackslash}p{(\linewidth - 6\tabcolsep) * \real{0.2308}}
  >{\raggedright\arraybackslash}p{(\linewidth - 6\tabcolsep) * \real{0.2308}}@{}}
\toprule\noalign{}
\begin{minipage}[b]{\linewidth}\raggedright
Componente nuevo
\end{minipage} & \begin{minipage}[b]{\linewidth}\raggedleft
Peso
\end{minipage} & \begin{minipage}[b]{\linewidth}\raggedright
¿Está ya dentro del Riesgo?
\end{minipage} & \begin{minipage}[b]{\linewidth}\raggedright
Veredicto
\end{minipage} \\
\midrule\noalign{}
\endhead
\bottomrule\noalign{}
\endlastfoot
Riesgo de incendios forestales & 30 \% & --- (es el contenedor) & --- \\
\textbf{Interfaz} & \textbf{20 \%} & \textbf{Sí, DOS veces}: como
Amenaza (§4.1.3.1) y como Vulnerabilidad (§5.1.5) & \textbf{Triple
conteo} \\
\textbf{Infraestructura crítica} & \textbf{20 \%} & \textbf{Sí,
parcialmente}: §5.1.4, pero sólo la energética & \textbf{Doble conteo
parcial} \\
Cicatrices de áreas afectadas & 15 \% & Solapamiento conceptual con la
Amenaza Estadística (§4.1.1), \textbf{pero en temporadas distintas} &
Riesgo bajo, con matices \\
Comunidades y Escuelas Preparadas & 15 \% & No aparece en el documento &
\textbf{Limpio} \\
\end{longtable}

El problema más serio no es el que parece. \textbf{La interfaz
urbano-forestal ya se cuenta dos veces dentro del propio índice de
Riesgo de CONAF}, antes de que nosotros le sumemos un 20 \% aparte. Al
sumarlo, queda contada tres veces.

\begin{center}\rule{0.5\linewidth}{0.5pt}\end{center}

\subsection{2. Cómo se construye el Riesgo, según el
documento}\label{cuxf3mo-se-construye-el-riesgo-seguxfan-el-documento}

El documento es explícito en su estructura (§4, §5, §5.2):

\begin{verbatim}
RIESGO = suma ponderada de  AMENAZA  +  VULNERABILIDAD        (pág. 26)

AMENAZA (pág. 9-18)
├── 4.1.1 Amenaza Estadística
│     └── Frecuencia de incendios, últimas 5 temporadas (2015-2016 a 2019-2020)
├── 4.1.2 Amenaza Estructural
│     ├── Velocidad de propagación por modelo de combustible
│     ├── Temperatura máxima (promedio histórico)
│     ├── Pendiente
│     └── Exposición
└── 4.1.3 Elementos de Amenaza  (antrópicos)
      ├── 4.1.3.1  INTERFAZ URBANO-FORESTAL        ←──┐
      ├── 4.1.3.2  Red de carreteras y caminos        │
      ├── 4.1.3.3  Densidad poblacional               │  la misma variable
      └── 4.1.3.4  Red eléctrica                      │  entra por los dos
                                                      │  lados del índice
VULNERABILIDAD (pág. 19-25)                           │
├── 5.1.1 Valor económico del combustible             │
├── 5.1.2 Valor ecológico del combustible             │
├── 5.1.3 SNASPE (áreas silvestres protegidas)        │
├── 5.1.4  INFRAESTRUCTURA CRÍTICA                    │
├── 5.1.5  INTERFAZ URBANO-FORESTAL              ←────┘
├── 5.1.6 Red eléctrica
└── 5.1.7 Red de caminos y carreteras
\end{verbatim}

Citas literales que sostienen lo anterior:

\begin{quote}
\textbf{§5.2, pág. 26} --- «La integración de los factores que se han
desarrollado con anterioridad, es decir, amenaza y vulnerabilidad, se
realiza mediante \textbf{la suma ponderada de los valores
--reescalados-- del territorio} para cada uno de ellos.»
\end{quote}

\begin{quote}
\textbf{§4.1.3.1, pág. 13} --- «Para la clasificación de las áreas de
\textbf{interfaz urbano forestal}, se han utilizado datos entregados por
el MINVU del Censo 2017. Se clasificó la información por tipo de área
urbana, resultando como área urbana consolidada principal, área urbana
consolidada secundaria, aldea y vivienda rural.»
\end{quote}

\begin{quote}
\textbf{§5.1.5, pág. 22} --- «La Figura 20 muestra el \textbf{área de
interfaz} para los diferentes tipos de viviendas de la región de Ñuble.
Las áreas urbanas consolidadas poseen una menor
\textbf{vulnerabilidad}\ldots»
\end{quote}

\begin{quote}
\textbf{§5.1.4, pág. 21} --- «La Figura 19 muestra las
\textbf{infraestructuras críticas de la superintendencia de electricidad
y combustible}\ldots{} principalmente \textbf{subestaciones eléctricas,
estaciones de servicios, centrales eléctricas, almacenamiento de
combustible} (gas, diesel).»
\end{quote}

Y el propio documento confirma, al describir los resultados, que la
interfaz es uno de los motores del resultado final:

\begin{quote}
\textbf{Cierre de §4.1, pág. 18} --- «Para la categoría muy alta y alta
la amenaza\ldots{} principalmente corresponde a las variables de
\textbf{interfaz urbano forestal} de las grandes ciudades\ldots{}
también las zonas que presentan una mayor frecuencia de incendios
históricos y una alta densidad de viviendas rurales.»

La misma página explica cómo se integran los tres componentes: «Para
obtener el factor de amenaza, \textbf{se intersectan} los índices
amenaza estadística, amenaza estructural y elementos de amenaza\ldots{}
después de realizar una agrupación mediante \textbf{cuantiles y criterio
experto}.»
\end{quote}

\begin{quote}
\textbf{§5.2 (resultados de riesgo), pág. 26} --- «Las zonas con riesgo
alto a muy alto\ldots{} coinciden con las áreas silvestres protegidas de
la región, incendios históricos, \textbf{interfaz urbano forestal} y
caminos que poseen alta transitabilidad.»
\end{quote}

\subsubsection{2.1 Recuento: cuántas veces entra cada
variable}\label{recuento-cuuxe1ntas-veces-entra-cada-variable}

Sumando lo que hace el documento (columnas 1 y 2) y lo que añaden las
ponderaciones nuevas (columna 3):

\begin{longtable}[]{@{}
  >{\raggedright\arraybackslash}p{(\linewidth - 8\tabcolsep) * \real{0.1579}}
  >{\centering\arraybackslash}p{(\linewidth - 8\tabcolsep) * \real{0.2105}}
  >{\centering\arraybackslash}p{(\linewidth - 8\tabcolsep) * \real{0.2105}}
  >{\centering\arraybackslash}p{(\linewidth - 8\tabcolsep) * \real{0.2105}}
  >{\centering\arraybackslash}p{(\linewidth - 8\tabcolsep) * \real{0.2105}}@{}}
\toprule\noalign{}
\begin{minipage}[b]{\linewidth}\raggedright
Variable
\end{minipage} & \begin{minipage}[b]{\linewidth}\centering
En Amenaza
\end{minipage} & \begin{minipage}[b]{\linewidth}\centering
En Vulnerab.
\end{minipage} & \begin{minipage}[b]{\linewidth}\centering
Componente aparte
\end{minipage} & \begin{minipage}[b]{\linewidth}\centering
\textbf{Total}
\end{minipage} \\
\midrule\noalign{}
\endhead
\bottomrule\noalign{}
\endlastfoot
\textbf{Interfaz urbano-forestal} & ✓ §4.1.3.1 & ✓ §5.1.5 & ✓ 20 \% &
\textbf{3×} \\
\textbf{Red eléctrica / subestaciones} & ✓ §4.1.3.4 & ✓ §5.1.6 + §5.1.4
& ✓ (dentro de infra) & \textbf{3-4×} \\
Red de caminos y carreteras & ✓ §4.1.3.2 & ✓ §5.1.7 & --- &
\textbf{2×} \\
Infraestructura crítica no eléctrica & --- & --- & ✓ 20 \% & 1× \\
Densidad poblacional & ✓ §4.1.3.3 & --- & --- & 1× \\
Frecuencia histórica de incendios & ✓ §4.1.1 & --- & --- & 1× \\
Combustible / clima / topografía & ✓ §4.1.2 & ✓ §5.1.1-5.1.2 & --- &
2× \\
SNASPE & --- & ✓ §5.1.3 & --- & 1× \\
Cicatrices 2021-2026 & --- & --- & ✓ 15 \% & 1× \\
Comunidades y Escuelas Preparadas & --- & --- & ✓ 15 \% & 1× \\
\end{longtable}

\textbf{Dos cosas que conviene separar al discutirlo:}

\begin{itemize}
\tightlist
\item
  \textbf{Lo que hereda del documento} (filas con ✓ en las dos primeras
  columnas): la interfaz, la red eléctrica, los caminos y el combustible
  ya entran dos veces \textbf{dentro del propio índice de CONAF}. Eso no
  lo introducimos nosotros y no lo podemos deshacer sin las capas
  intermedias. Es discutible que sea un error: en el marco amenaza ×
  vulnerabilidad, una misma variable puede legítimamente actuar como
  fuente de ignición \emph{y} como bien expuesto. Lo que no es
  defendible es \textbf{no declararlo}.
\item
  \textbf{Lo que añadimos nosotros} (columna 3): aquí sí hay decisión
  propia, y es donde la interfaz pasa de 2× a 3×.
\end{itemize}

\begin{center}\rule{0.5\linewidth}{0.5pt}\end{center}

\subsection{3. Caso por caso}\label{caso-por-caso}

\subsubsection{\texorpdfstring{3.1 Interfaz --- triple conteo.
\textbf{El problema más
grave.}}{3.1 Interfaz --- triple conteo. El problema más grave.}}\label{interfaz-triple-conteo.-el-problema-muxe1s-grave.}

La interfaz entra \textbf{tres veces} en el puntaje final:

\begin{enumerate}
\def\labelenumi{\arabic{enumi}.}
\tightlist
\item
  Dentro del 30 \% de Riesgo, vía \textbf{Amenaza} (§4.1.3.1): la
  interfaz como fuente de ignición --- viviendas rurales, uso cotidiano
  del fuego, quemas.
\item
  Dentro del mismo 30 \% de Riesgo, vía \textbf{Vulnerabilidad}
  (§5.1.5): la interfaz como elemento expuesto a dañarse.
\item
  Como componente propio, con \textbf{20 \%}.
\end{enumerate}

Las dos primeras están dentro del documento de CONAF y no dependen de
nosotros. La tercera es la que añade la ponderación nueva.

\textbf{Que es la misma capa, no una parecida.} El documento clasifica
la interfaz en «área urbana consolidada principal, área urbana
consolidada secundaria, aldea y vivienda rural» (MINVU, Censo 2017). La
capa \texttt{raw2/INTERFAZ} que alimenta nuestro componente tiene
exactamente esas categorías, más dos:

\begin{verbatim}
Vivienda Rural          340.877
Urbano Principal         12.481
Urbano Secundario        11.541
Loteos - Parcelaciones    6.550
Aldea                     4.276
Campamentos               1.216
                        ────────
                        376.941 polígonos
\end{verbatim}

Las cuatro categorías del documento están ahí con el mismo nombre. No es
una fuente alternativa: es la misma familia de datos.

\subsubsection{\texorpdfstring{3.2 Infraestructura crítica --- doble
conteo
\textbf{parcial}}{3.2 Infraestructura crítica --- doble conteo parcial}}\label{infraestructura-cruxedtica-doble-conteo-parcial}

Aquí conviene no exagerar. La infraestructura del documento es
\textbf{energética y de combustibles}, de la SEC: subestaciones
eléctricas, estaciones de servicio, centrales eléctricas, almacenamiento
de gas y diésel.

La nuestra (\texttt{raw2/INFRAESTRUCTURA\_CRITICA}) tiene seis familias:

\begin{longtable}[]{@{}
  >{\raggedright\arraybackslash}p{(\linewidth - 2\tabcolsep) * \real{0.5000}}
  >{\raggedright\arraybackslash}p{(\linewidth - 2\tabcolsep) * \real{0.5000}}@{}}
\toprule\noalign{}
\begin{minipage}[b]{\linewidth}\raggedright
Familia
\end{minipage} & \begin{minipage}[b]{\linewidth}\raggedright
¿En el documento CONAF?
\end{minipage} \\
\midrule\noalign{}
\endhead
\bottomrule\noalign{}
\endlastfoot
SUBESTACIONES & \textbf{Sí} (§5.1.4, y además §4.1.3.4 y §5.1.6 como red
eléctrica) \\
ANTENAS & No \\
CENTROS\_DE\_SALUD & No \\
CENTROS\_PENITENCIARIOS & No \\
RED\_AEROPORTUARIA & No \\
SERVICIOS\_SANITARIOS\_RURALES & No \\
\end{longtable}

\textbf{Sólo una de las seis familias se duplica.} Las otras cinco son
aportes genuinamente nuevos: el índice de CONAF no valora hospitales,
escuelas, antenas ni cárceles.

Ahora bien, el solapamiento de la parte eléctrica es \textbf{triple},
igual que la interfaz: la red eléctrica está en Amenaza (§4.1.3.4) y en
Vulnerabilidad (§5.1.6), y las subestaciones otra vez en Infraestructura
Crítica (§5.1.4).

\textbf{Matiz sobre el peso real.} En el modelo actual las siete
familias se promedian con peso igual --- verificado en
\texttt{priorizacion.py}, \texttt{pesos\_infra} asigna \texttt{1.0} a
subestaciones, antenas, aeroportuaria, ssr, salud, educacion y
penitenciarias ---, así que las subestaciones aportan \textbf{1/7 del
subíndice}. Contando puntos publicados son \textbf{26 de 696 (3,7 \%)}.
El doble conteo existe, pero su magnitud dentro de este componente es
pequeña.

(Nota: el modelo actual usa siete familias porque incluye
\texttt{educacion}; la carpeta \texttt{raw2/INFRAESTRUCTURA\_CRITICA} de
la metodología nueva trae seis, sin establecimientos educacionales.
Conviene confirmar si eso es deliberado.)

\subsubsection{\texorpdfstring{3.3 Cicatrices --- el solapamiento que
parece obvio y \textbf{no lo
es}}{3.3 Cicatrices --- el solapamiento que parece obvio y no lo es}}\label{cicatrices-el-solapamiento-que-parece-obvio-y-no-lo-es}

La sospecha natural es que las cicatrices dupliquen la \textbf{Amenaza
Estadística}, que es «el índice frecuencia; relativo al número de
incendios forestales de las últimas 5 temporadas» (§4.1.1, pág. 9). Pero
hay tres diferencias reales:

\begin{enumerate}
\def\labelenumi{\arabic{enumi}.}
\tightlist
\item
  \textbf{Las ventanas temporales no se solapan.} El documento usa
  \textbf{2015-2016 a 2019-2020}. Las cicatrices disponibles
  (\texttt{raw2/POLIGONO\_DE\_AREA\_AFECTADA}) son \textbf{2021-2022,
  2022-2023, 2023-2024, 2024-2025 y 2025-2026}. Son periodos
  \textbf{disjuntos}: el índice de Riesgo de 2021 no puede contener
  información de incendios posteriores a su publicación.
\item
  \textbf{Miden cosas distintas.} Frecuencia = número de eventos.
  Cicatriz = superficie efectivamente quemada. Un sector con muchos
  incendios pequeños y otro con un solo megaincendio pueden intercambiar
  posiciones según cuál se mire.
\item
  \textbf{El signo puede ser opuesto.} Una cicatriz reciente
  \textbf{reduce} el combustible disponible. El índice de Riesgo ya
  pondera «velocidad de propagación de los modelos de combustible»
  (§4.1.2.1), que en un área recién quemada es baja. Sumar la cicatriz
  como factor que \textbf{aumenta} la prioridad puede estar
  contradiciendo, no duplicando, lo que ya dice el Riesgo.
\end{enumerate}

\textbf{Conclusión:} no es doble conteo en sentido estricto. Pero el
punto 3 merece una decisión explícita del equipo: \emph{¿una cicatriz
reciente hace un sector más prioritario o menos?} Las dos respuestas son
defendibles (más: el sector demostró ser propenso y hay que recuperarlo;
menos: no hay combustible que arda en los próximos años), pero hay que
elegir una y escribirla, porque el modelo hoy no la tiene.

\subsubsection{3.4 Comunidades y Escuelas Preparadas ---
limpio}\label{comunidades-y-escuelas-preparadas-limpio}

No aparece en ninguna sección del documento. Es un dato de gestión de
CONAF (dónde se ha intervenido), no una característica del territorio.
Es el único de los cuatro componentes externos que no arrastra ningún
solapamiento.

Ojo con una cosa distinta: «Preparadas» significa que \textbf{ya se
trabajó ahí}. Si el índice lo suma como factor que aumenta la prioridad,
está priorizando lo ya atendido. El modelo actual tiene esto resuelto a
medias --- calcula un \texttt{indice\_brecha} que invierte el signo de
este componente (\texttt{1.0\ -\ sub\_preparadas}) para responder
«¿dónde todavía no hemos llegado?» --- pero \textbf{ese índice no es el
que se publica}. Es otra decisión que conviene hacer explícita.

\begin{center}\rule{0.5\linewidth}{0.5pt}\end{center}

\subsection{4. Evidencia empírica sobre los datos
publicados}\label{evidencia-empuxedrica-sobre-los-datos-publicados}

Correlación de Spearman entre los subíndices de las 572 manchas de
\texttt{priorizacion.geojson} (Coyhaique, Los Ángeles y Mulchén):

\begin{longtable}[]{@{}lr@{}}
\toprule\noalign{}
Par & ρ \\
\midrule\noalign{}
\endhead
\bottomrule\noalign{}
\endlastfoot
sub\_riesgo --- sub\_interfaz & \textbf{+0,486} \\
sub\_interfaz --- sub\_infra & +0,461 \\
sub\_riesgo --- sub\_infra & +0,329 \\
sub\_riesgo --- sub\_preparadas & +0,171 \\
sub\_interfaz --- sub\_preparadas & +0,189 \\
sub\_infra --- sub\_preparadas & +0,202 \\
\end{longtable}

\textbf{Cómo hay que leer esto, con honestidad.} Una correlación de
+0,486 entre riesgo e interfaz es \textbf{consistente} con que la
interfaz esté dentro del riesgo, pero \textbf{no lo demuestra}: podría
explicarse igual por el hecho trivial de que donde hay gente hay
interfaz, infraestructura y también más ignición. \textbf{La prueba del
doble conteo es documental} (sección 2), no estadística. La correlación
sólo confirma que el efecto no es despreciable en los datos reales.

Sí es concluyente lo otro: los dos componentes que se duplican son
justamente los que más territorio cubren.

\begin{longtable}[]{@{}lrr@{}}
\toprule\noalign{}
Subíndice & Manchas con valor \textgreater{} 0 & Media \\
\midrule\noalign{}
\endhead
\bottomrule\noalign{}
\endlastfoot
sub\_riesgo & 570 / 572 (99,7 \%) & 0,5268 \\
sub\_interfaz & 436 / 572 (76,2 \%) & 0,3597 \\
sub\_infra & 108 / 572 (18,9 \%) & 0,0259 \\
sub\_preparadas & 33 / 572 (5,8 \%) & 0,0184 \\
\end{longtable}

Riesgo e interfaz están presentes en casi todas partes; infraestructura
y preparadas son casi siempre cero. \textbf{El solapamiento se concentra
exactamente en los dos componentes que deciden el mapa.}

\begin{center}\rule{0.5\linewidth}{0.5pt}\end{center}

\subsection{5. Qué problemas causa esto para la
metodología}\label{quuxe9-problemas-causa-esto-para-la-metodologuxeda}

\subsubsection{5.1 Los pesos declarados dejan de significar lo que
dicen}\label{los-pesos-declarados-dejan-de-significar-lo-que-dicen}

Es el problema central. Cuando el programa declara «Interfaz: 20 \%»,
cualquiera entiende que la interfaz determina una quinta parte del
resultado. No es así.

Sea \textbf{α} la fracción del índice de Riesgo atribuible a la interfaz
(por sus dos entradas, amenaza y vulnerabilidad). El peso efectivo es:

\begin{verbatim}
peso_efectivo(interfaz) = 0,20 + 0,30 · α
\end{verbatim}

\begin{longtable}[]{@{}lrr@{}}
\toprule\noalign{}
Si α vale\ldots{} & Peso efectivo & Sobre el declarado \\
\midrule\noalign{}
\endhead
\bottomrule\noalign{}
\endlastfoot
0,10 & 0,230 & +15 \% \\
0,20 & 0,260 & +30 \% \\
0,30 & 0,290 & +45 \% \\
0,40 & 0,320 & +60 \% \\
\end{longtable}

\textbf{No podemos fijar α}, porque el documento \textbf{no publica los
pesos internos} de la suma ponderada (ver §7, incertidumbres). Pero
cualquier valor razonable deja el peso efectivo de la interfaz por
encima del de infraestructura, que nominalmente es igual (20 \%). Los
dos componentes que el programa declara equivalentes \textbf{no lo son}.

\subsubsection{5.2 El análisis de sensibilidad da resultados
falsos}\label{el-anuxe1lisis-de-sensibilidad-da-resultados-falsos}

Si alguien pregunta «¿cuánto cambia el mapa si bajamos la interfaz del
20 \% al 10 \%?», la respuesta que se obtenga será \textbf{una
subestimación}: al mover ese peso sólo se mueve una de las tres entradas
de la interfaz; las otras dos viajan dentro del 30 \% de Riesgo y no se
tocan. El modelo parecerá más robusto a ese parámetro de lo que
realmente es.

Esto no es hipotético: el notebook ya mide que, moviendo pesos ±0,10,
cambian de clase entre el 21,8 \% y el 28,1 \% de los hexágonos. Esa
cifra está calculada sobre pesos nominales, así que describe menos
movimiento del que habría.

\subsubsection{5.3 Doble penalización territorial y refuerzo
circular}\label{doble-penalizaciuxf3n-territorial-y-refuerzo-circular}

Un sector de interfaz urbano-forestal con una subestación recibe puntaje
por:

\begin{itemize}
\tightlist
\item
  estar en interfaz (20 \% directo),
\item
  estar clasificado como riesgo alto \textbf{porque tiene interfaz}
  (dentro del 30 \%),
\item
  tener infraestructura crítica (20 \%),
\item
  estar clasificado como riesgo alto \textbf{porque tiene red eléctrica}
  (otra vez dentro del 30 \%).
\end{itemize}

Los factores no se suman, se \textbf{refuerzan}. El resultado es que las
zonas de interfaz saldrán sistemáticamente como las más prioritarias.
Puede que eso sea correcto como política ---proteger donde vive la
gente---, pero deja de ser \textbf{una conclusión del análisis} para ser
\textbf{un artefacto de su construcción}. Y esa diferencia es justamente
la que un análisis territorial existe para aportar.

\subsubsection{5.4 Pérdida de poder
discriminante}\label{puxe9rdida-de-poder-discriminante}

Cuando varios componentes miden lo mismo, el índice compuesto se
aproxima a una función de un solo factor y deja de separar territorios.
En los datos publicados eso todavía no es agudo (ρ = +0,486 entre riesgo
e interfaz, sección 4), pero el síntoma ya apareció antes en este
proyecto: el notebook dejó registrada una correlación de \textbf{0,85
entre \texttt{sub\_interfaz} y \texttt{sub\_infra} en la capa gruesa} de
40.000 ha, con aviso explícito de redundancia. \emph{(Ese 0,85 procede
del notebook y no lo recalculé: la capa gruesa tiene 15 manchas y no
está publicada. Las cifras de la sección 4 sí son medición propia sobre
las 572 manchas.)}

\subsubsection{5.5 El peso nominal ya no se parece al efectivo, y está
medido}\label{el-peso-nominal-ya-no-se-parece-al-efectivo-y-estuxe1-medido}

Con las ponderaciones \textbf{actuales} (riesgo 0,40 · interfaz 0,20 ·
infra 0,20 · preparadas 0,20), el peso \textbf{efectivo} de cada
componente ---su aporte real a \texttt{puntaje\_medio}--- no se parece
al nominal. Calculado sobre las 572 manchas publicadas:

\begin{longtable}[]{@{}lrrr@{}}
\toprule\noalign{}
Componente & Peso nominal & Aporte medio & \textbf{Peso efectivo} \\
\midrule\noalign{}
\endhead
\bottomrule\noalign{}
\endlastfoot
Riesgo & 0,40 & 0,2107 & \textbf{0,723} \\
Interfaz & 0,20 & 0,0719 & \textbf{0,247} \\
Infraestructura & 0,20 & 0,0052 & \textbf{0,018} \\
Preparadas & 0,20 & 0,0037 & \textbf{0,013} \\
\end{longtable}

\emph{(Control aritmético: los cuatro aportes suman 0,2915, que es
exactamente la media observada de \texttt{puntaje\_medio}.)}

La causa aquí es distinta del doble conteo ---infra y preparadas valen
cero en el 81 \% y el 94 \% de las manchas, así que no pueden mover
nada---, pero el efecto se suma al anterior: \textbf{el mapa es, de
hecho, riesgo (72 \%) + interfaz (25 \%); los otros dos componentes
juntos no llegan al 4 \%} pese a declarar 20 \% cada uno. Cambiar los
pesos nominales de 40/20/20/20 a 30/20/20/15/15 no corregirá eso por sí
solo: mientras un componente sea cero en casi todo el territorio, su
peso nominal es decorativo.

Y nótese la coincidencia incómoda: \textbf{los dos componentes que sí
mueven el mapa (riesgo e interfaz, 97 \% del aporte) son exactamente los
dos que se solapan entre sí.}

\subsubsection{5.6 Un agravante práctico: el Riesgo llega como caja
negra}\label{un-agravante-pruxe1ctico-el-riesgo-llega-como-caja-negra}

Comprobado sobre el insumo real: la capa \texttt{RIESGO\_3\_COMUNAS.shp}
(463.366 polígonos) tiene \textbf{exactamente dos campos}: \texttt{Id} y
\texttt{gridcode} (0 a 4, de Muy Bajo a Muy Alto).

\begin{verbatim}
campos: ['Id', 'gridcode']
gridcode: {2: 20082, 3: 17392, 1: 10737, 4: 1416, 0: 373}
\end{verbatim}

Es el producto \textbf{final} ya integrado. No trae la amenaza y la
vulnerabilidad por separado, ni ninguno de sus subcomponentes.
\textbf{Con este insumo es imposible descontar la contribución de la
interfaz o de la red eléctrica}: no se puede separar lo que ya viene
sumado. Esto descarta la solución técnicamente más limpia (reconstruir
el riesgo sin los componentes repetidos) salvo que GEPRIF entregue las
capas intermedias.

\begin{center}\rule{0.5\linewidth}{0.5pt}\end{center}

\subsection{6. Opciones, con su coste}\label{opciones-con-su-coste}

Ninguna es gratis. Van de menor a mayor esfuerzo.

\subsubsection{Opción A --- Declararlo y no tocar los
pesos}\label{opciuxf3n-a-declararlo-y-no-tocar-los-pesos}

Mantener 30/20/20/15/15 y documentar que \textbf{son pesos de énfasis,
no de contribución}, con la tabla de la sección 3 al lado.

\begin{itemize}
\tightlist
\item
  \textbf{A favor:} coste cero; no rompe la comparabilidad con lo ya
  entregado; políticamente es defendible que la interfaz pese más,
  porque es donde está la gente.
\item
  \textbf{En contra:} el índice sigue sin poder interpretarse como se
  lee, y el análisis de sensibilidad sigue mintiendo. Traslada el
  problema al lector.
\end{itemize}

\subsubsection{Opción B --- Quitar del cálculo lo que ya está dentro del
Riesgo}\label{opciuxf3n-b-quitar-del-cuxe1lculo-lo-que-ya-estuxe1-dentro-del-riesgo}

Eliminar «Interfaz» como componente propio y dejar que entre sólo por el
Riesgo; dejar en «Infraestructura crítica» únicamente las cinco familias
que CONAF no considera (antenas, salud, penitenciarios, aeroportuaria,
SSR).

\begin{itemize}
\tightlist
\item
  \textbf{A favor:} elimina el doble conteo de raíz; cada elemento entra
  una vez.
\item
  \textbf{En contra:} se pierde control sobre cuánto pesa la interfaz,
  porque queda sepultada en un peso interno que no conocemos; y el
  resultado cambiará bastante respecto de lo ya mostrado. Además, el 20
  \% liberado hay que repartirlo, y eso es otra decisión.
\end{itemize}

\subsubsection{Opción C --- Pedir a GEPRIF las capas
intermedias}\label{opciuxf3n-c-pedir-a-geprif-las-capas-intermedias}

Solicitar al Departamento de Desarrollo e Investigación la Amenaza y la
Vulnerabilidad por separado, o mejor, los subíndices. Con ellos se puede
construir un índice \textbf{sin repeticiones}, eligiendo explícitamente
por dónde entra cada variable.

\begin{itemize}
\tightlist
\item
  \textbf{A favor:} es la única solución metodológicamente correcta;
  además permitiría publicar los pesos reales.
\item
  \textbf{En contra:} depende de terceros y de plazos que no
  controlamos. Es la vía a abrir ya, aunque se adopte A o B mientras
  tanto.
\end{itemize}

\subsubsection{Opción D --- Sustituir el Riesgo por sus partes no
repetidas}\label{opciuxf3n-d-sustituir-el-riesgo-por-sus-partes-no-repetidas}

Usar el Riesgo \textbf{sólo} por su Amenaza Estadística y su Amenaza
Estructural (frecuencia, combustible, temperatura, pendiente,
exposición) y llevar interfaz, infraestructura y red vial fuera, como
componentes propios y con peso explícito.

\begin{itemize}
\tightlist
\item
  \textbf{A favor:} es la formulación más limpia y la más fácil de
  explicar.
\item
  \textbf{En contra:} requiere lo mismo que C (capas intermedias) y
  además rehacer la metodología, no sólo los pesos.
\end{itemize}

\subsubsection{Recomendación}\label{recomendaciuxf3n}

\textbf{Abrir la Opción C ya} ---cuesta un correo y es la única que
resuelve el fondo--- y mientras llega la respuesta \textbf{adoptar la A
con la tabla de solapamientos publicada junto al mapa}. La B es
tentadora por limpia, pero cambiar el resultado sin poder explicar
cuánto pesa ahora la interfaz sustituye un problema conocido por otro
opaco.

Con independencia de la opción, hay \textbf{dos decisiones que el equipo
debe tomar y escribir}, porque hoy no están en ninguna parte y no son
técnicas:

\begin{enumerate}
\def\labelenumi{\arabic{enumi}.}
\tightlist
\item
  \textbf{¿Una cicatriz reciente sube o baja la prioridad?} (sección
  3.3)
\item
  \textbf{¿«Preparadas» suma o resta?} ¿El mapa señala dónde hay que
  trabajar, o dónde ya se trabajó? (sección 3.4)
\end{enumerate}

\begin{center}\rule{0.5\linewidth}{0.5pt}\end{center}

\subsection{7. Incertidumbres --- lo que NO se pudo
comprobar}\label{incertidumbres-lo-que-no-se-pudo-comprobar}

Se declaran porque afectan a la fuerza de las conclusiones.

\begin{enumerate}
\def\labelenumi{\arabic{enumi}.}
\tightlist
\item
  \textbf{No conozco los pesos internos del índice de Riesgo.} El
  documento dice «suma ponderada» (§5.2) pero \textbf{no publica los
  coeficientes} de amenaza ni de vulnerabilidad, ni el reparto entre sus
  subcomponentes. Por eso la sección 5.1 está planteada de forma
  paramétrica (en función de α) y no con un número. \textbf{Sin ese dato
  no se puede cuantificar el sobrepeso exacto, sólo acotarlo.}
\item
  \textbf{Tampoco se conoce cómo se «intersectan» los tres componentes
  de la amenaza.} El documento dice que se intersectan y luego se
  agrupan «mediante cuantiles y criterio experto» (pág. 18). «Criterio
  experto» no es reproducible desde el documento.
\item
  \textbf{No verifiqué el solapamiento espacial píxel a píxel} entre la
  capa de riesgo alto y la de interfaz. Habría sido la evidencia más
  directa, pero con \texttt{gridcode} como único atributo (§5.6) sólo
  mediría coincidencia territorial, que no distingue «la interfaz está
  dentro del riesgo» de «ambos ocurren donde hay gente». Las
  correlaciones de la sección 4 se aportan con esa misma cautela.
\item
  \textbf{Las ponderaciones nuevas no están implementadas todavía.} El
  modelo publicado sigue con 0,40 / 0,20 / 0,20 / 0,20 y \textbf{sin
  componente de cicatrices}. Todo lo anterior es un análisis previo a
  implementarlas.
\item
  \textbf{La capa de cicatrices 2025-2026 no está completa}; la carpeta
  se llama literalmente \texttt{2025-2026\ (se\ cargará\ lunes\ 14)}.
\end{enumerate}

\begin{center}\rule{0.5\linewidth}{0.5pt}\end{center}

\subsection{8. Resumen para la
reunión}\label{resumen-para-la-reuniuxf3n}

\begin{itemize}
\tightlist
\item
  \textbf{Sí hay doble ponderación, y el caso grave es la interfaz}: ya
  entra dos veces dentro del índice de Riesgo de CONAF (como amenaza y
  como vulnerabilidad), así que sumarle un 20 \% propio la cuenta
  \textbf{tres veces}.
\item
  \textbf{Infraestructura crítica se duplica sólo en su parte eléctrica}
  (subestaciones, 1 de 6 familias). Las otras cinco son aporte genuino.
\item
  \textbf{Cicatrices no es doble conteo}: cubre temporadas
  \emph{posteriores} (2021-2026) a las del índice de Riesgo (2015-2020).
  Pero hay que decidir si una cicatriz reciente sube o baja la
  prioridad, porque el combustible quemado no vuelve a arder de
  inmediato.
\item
  \textbf{Preparadas está limpio}, aunque hay que decidir si suma o
  resta.
\item
  \textbf{Consecuencia práctica:} los pesos declarados no son los
  efectivos ---hoy el mapa es riesgo 72 \% + interfaz 25 \%, con
  infraestructura y preparadas por debajo del 4 \% juntas, pese a
  declarar 20 \% cada una---, el análisis de sensibilidad subestima el
  papel de la interfaz, y las zonas de interfaz saldrán siempre arriba
  por construcción y no por hallazgo.
\item
  \textbf{Obstáculo:} el Riesgo llega como un único \texttt{gridcode} de
  0 a 4, así que hoy \textbf{no se puede descontar} lo repetido. La
  salida de fondo es pedir a GEPRIF las capas intermedias.
\end{itemize}

\end{document}
